\documentclass[10pt,twocolumn,a4paper]{article}

\usepackage[a4paper,margin=1.9cm,top=1.8cm,bottom=2.0cm,columnsep=0.7cm]{geometry}
\usepackage{fontspec}
\setmainfont{CMU Serif}
\setmonofont{CMU Typewriter Text}
\usepackage[main=ukrainian,english]{babel}
\usepackage{microtype}
\usepackage{amsmath,amssymb,amsfonts}
\usepackage{graphicx}
\usepackage{booktabs}
\usepackage{array}
\usepackage{float}
\usepackage{caption}
\usepackage{subcaption}
\usepackage{enumitem}
\usepackage{titlesec}
\usepackage{xcolor}
\usepackage{hyperref}
\usepackage{fancyhdr}
\usepackage{balance}

\hypersetup{
  colorlinks=true,
  linkcolor=blue!40!black,
  citecolor=blue!40!black,
  urlcolor=blue!40!black,
  pdftitle={Аналіз аудіофіч треків Spotify та прогноз популярності},
  pdfauthor={Масич Нікіта}
}

\pagestyle{fancy}
\fancyhf{}
\fancyhead[L]{Інтелектуальний аналіз даних}
\fancyhead[R]{\thepage}
\renewcommand{\headrulewidth}{0.4pt}
\renewcommand{\footrulewidth}{0pt}

\setlength{\parindent}{0pt}
\setlength{\parskip}{0.5em}
\setlength{\columnsep}{0.7cm}
\setlist[itemize]{leftmargin=*, itemsep=0.2em, topsep=0.2em}
\setlist[enumerate]{leftmargin=*, itemsep=0.2em, topsep=0.2em}

\titleformat{\section}
{\large\bfseries}
{\thesection.}{0.5em}{}

\titleformat{\subsection}
{\normalsize\bfseries}
{\thesubsection.}{0.5em}{}

\renewcommand{\abstractname}{Анотація}

\newcommand{\ProjectTopic}{Аналіз аудіофіч треків Spotify та прогноз популярності}
\newcommand{\DisciplineName}{Інтелектуальний аналіз даних}
\newcommand{\StudentName}{Масич Нікіта}
\newcommand{\StudentGroup}{АСД}
\newcommand{\UniversityName}{Київський національний університет імені Тараса Шевченка}
\newcommand{\FacultyName}{Факультет комп'ютерних наук та кібернетики}
\newcommand{\AcademicYear}{2026}

\newcommand{\SectionFile}[2]{%
  \IfFileExists{sections/#1}{%
    \input{sections/#1}%
  }{%
    \section{#2}
    \noindent Текст цієї секції буде додано пізніше в окремому файлі \texttt{sections/\detokenize{#1}}.%
  }%
}

\begin{document}

\twocolumn[
  \begin{center}
    {\Large\bfseries \ProjectTopic\par}
    \vspace{0.5em}
    {\normalsize Проєкт з дисципліни ``\DisciplineName''\par}
    \vspace{1em}
    \begin{tabular}{@{}c@{}}
      \textbf{Студент:} \StudentName \\
      \textbf{Група:} \StudentGroup \\
      \textbf{Курс:} 1 курс магістратури \\
      \textbf{Університет:} \UniversityName \\
      \textbf{Факультет:} \FacultyName \\
      \textbf{Навчальний рік:} \AcademicYear
    \end{tabular}
    \vspace{1em}
  \end{center}
]

\begin{abstract}
  У цій роботі досліджено набір аудіофіч треків Spotify з метою аналізу взаємозв'язків між характеристиками композицій, жанровими групами та популярністю треків. Робота поєднує описовий аналіз даних, візуалізацію, кластеризацію та базове машинне навчання для побудови інтерпретованого та відтворюваного дослідження.
\end{abstract}

\textbf{Ключові слова:} інтелектуальний аналіз даних, Spotify, кластеризація, класифікація, аудіофічі, популярність треків.

\SectionFile{01_intro.tex}{Вступ і мотивація}
\SectionFile{02_related_work.tex}{Поточний стан досліджень}
\SectionFile{03_methods_results.tex}{Методи та результати}
\SectionFile{04_conclusions.tex}{Висновки та подальші кроки}

\balance
\begin{thebibliography}{99}

  \bibitem{spotify}
  Spotify Tracks Dataset. Kaggle dataset.

  \bibitem{scikit}
  Scikit-learn: Machine Learning in Python. User Guide.

  \bibitem{pandas}
  The pandas development team. Pandas documentation.

\end{thebibliography}

\end{document}
